\begin{textsample}{POS Dim 2 – al – Score 18.00 – al/al\_em\_46.txt}  \label{ex:f2_pos_270}
Para \textbf{itinerários} de Formação \textbf{Técnica} e Profissional, que apresentam 1200 horas no \textbf{curso} \textbf{técnico}, observa -se que Aprofundamento, \textbf{Eletivas} e Projetos de Vida contribuem na formação do \textbf{técnico}, portanto compõem as unidades curriculares da proposta curricular do \textbf{curso} \textbf{técnico}.
Contudo, neste modelo, a Unidade de Ensino não tem a flexibilidade demonstrada nos modelos com \textbf{cursos} \textbf{técnicos} de 800 e 1000 horas, precisando, por meio das propostas das unidades curriculares do \textbf{curso} \textbf{técnico} de 1200 horas, avançar na articulação com a formação geral, integração curricular e ações pedagógicas do território.
Observamos que, nesta distribuição da \textbf{carga} \textbf{horária} do IF, todas as unidades curriculares do aprofundamento, da \textbf{oferta} \textbf{eletiva} e do projeto de vida \textbf{concorrem} para a formação do \textbf{técnico} com \textbf{carga} \textbf{horária} de 1200 horas.
Os \textbf{itinerários} \textbf{formativos} de formação \textbf{técnica} e \textbf{profissional} devem garantir uma Formação para o Mundo do Trabalho.
Para isto, é preciso que os projetos apresentados na \textbf{trajetória} do EM garantam também propostas de Práticas \textbf{Profissionais}, enquanto unidades do \textbf{itinerário} de EPT, como os estágios supervisionados obrigatórios ou não, bem como programas de aprendizagem ligados às empresas, instituições e território local, regional e nacional.
É importante destacar que os estágios poderão fazer parte da proposta curricular nos casos em que a escolha da unidade de ensino esteja intimamente vinculada ao \textbf{curso} em que haja a sua obrigatoriedade pela instituição da classe credenciada, como, por exemplo, os \textbf{cursos} da área de saúde que atendem diretamente às unidades hospitalares ou setores de serviços afins.
A Lei nº 13.415/2017, em seu art. 4º, parágrafo 6º, oportuniza que, a critério dos sistemas de ensino, a \textbf{oferta} de formação com ênfase \textbf{técnica} e \textbf{profissional} considere a inclusão de vivências práticas de trabalho no setor produtivo ou em ambientes de simulação, bem como a possibilidade de concessão de certificados de \textbf{qualificação} para o trabalho, quando a formação for estruturada e organizada em etapas com terminalidade.
As unidades de ensino, conforme a proposição do \textbf{itinerário} \textbf{formativo} de formação \textbf{técnica} e \textbf{profissional}, devem apresentar dispositivos claros para o estudante, conforme \textbf{curso} \textbf{técnico}, \textbf{carga} \textbf{horária}, distribuição de unidades curriculares e tempos, instrumentos para matrícula, escolha e reposição desta formação.
Isto porque o estudante deve ter o conhecimento do tempo em que cada ano será apresentado em sua estrutura, a plena necessidade do cumprimento da \textbf{carga} \textbf{horária} e de seus conteúdos contidos nos eixos \textbf{estruturantes} que perpassam por sua formação.
Os \textbf{cursos} \textbf{técnicos} propostos no \textbf{itinerário} \textbf{formativo} serão acompanhados pela Inspeção dos Sistemas de Ensino, bem como poderão ser submetidos a avaliações de larga escala pela unidade de ensino ou pelo próprio sistema de ensino.
Conforme a proposta, o \textbf{itinerário} \textbf{formativo} de formação \textbf{técnica} e \textbf{profissional}, de acordo com critérios operacionais e legais, pode haver \textbf{certificação} ao final da educação básica ou intermediária com a conclusão de \textbf{cursos} FIC de \textbf{qualificação} \textbf{profissional}.
O ReCAL do EM tem como um dos grandes motivadores promover a melhor aderência e \textbf{mobilidade} entre as \textbf{ofertas} dos \textbf{itinerários} do sistema \textbf{estadual} de ensino de Alagoas, bem como as \textbf{ofertas} de EM pelo Brasil.
Neste sentido, recomenda -se para elaboração do IF de ETP : - Apresentação dos princípios norteadores para a (re)elaboração dos Planos de \textbf{Curso} dos \textbf{Cursos} \textbf{Técnicos}, dos \textbf{Cursos} de \textbf{Qualificação} \textbf{Profissional} ( referência Classificação Brasileira de Ocupações – CBO ) e do Programa de Aprendizagem, considerando : as competências para o mundo do trabalho ; o \textbf{perfil} do egresso de cada \textbf{curso} ; que as competências do \textbf{curso} técnico/qualificação \textbf{profissional} ou programa de aprendizagem reflitam os objetivos dos \textbf{cursos} ; e que as habilidades estejam atreladas às competências. - Orientações para que as matrizes curriculares para os \textbf{Cursos} \textbf{Técnicos} sejam elaboradas a partir das competências, articulando as diferentes habilidades por componentes curriculares. - Orientações para elaboração das ementas dos \textbf{Cursos} \textbf{Técnicos}, \textbf{Cursos} de \textbf{Qualificação} Profissional e Programa de Aprendizagem, que devem conter : título ; o que é ; por que é importante na formação do jovem ; articulação com o mundo do trabalho ; aprendizagens em jogo.
Indicação de que os Planos de \textbf{Curso} virão como anexo. - Eixos \textbf{estruturantes} : apresentação dos princípios para a articulação dos quatro eixos considerando as competências gerais da BNCC, a formação para o mundo do trabalho e as competências específicas da formação escolhida. - Orientações para que a articulação dos eixos \textbf{estruturantes} e suas competências específicas sejam em um módulo ou em componentes na parte de aprofundamento. - Orientações para a construção de estratégias didáticas que garantam as aprendizagens específicas dos eixos \textbf{estruturantes}. - \textbf{Certificação} : indicação de tipos de \textbf{certificação} ( intermediária, diploma ) e de quem certifica ( escola, parceiro etc. ).
Por fim, a Seduc irá estabelecer um cronograma para a implementação das alterações promovidas pela Lei 13.415/2017 para a Rede Pública Estadual, bem como a implementação deste ReCAL para a etapa Ensino Médio.
Apresentaremos, a seguir, matrizes ou modelos que apoiaram as \textbf{redes} de ensino na elaboração e planejamento da implementação da BNCC e IF.

% matched lemmas: carga, certificação, concorrer, curso, eletivo, estadual, estruturante, formativo, horário, itinerário, mobilidade, oferta, perfil, profissional, qualificação, rede, trajetória, técnico
\end{textsample}
